\chapter{Transistores Bipolares (BJT)}
\section{Generalidades}
???
\subsection{Símbolos}
???
\subsection{Corrientes y tensiones}
???
\section{Transistores NPN}
???
\subsection{Modelos de operación}
???
\subsection{Operación de un transistor NPN en modo normal}
???
\subsubsection{Curvas características}
???
\subsubsection{Relaciones teóricas}
???
\subsection{Operación de un transistor en continua}
???
\subsubsection{Solución gráfica}
???
\subsubsection{Modelos del transistor de corriente continua}
???
\subsubsection{Puntos de operación, punto de reposo o punto Q}
???
\subsection{Circuitos de polarización en región negativa}
???
\subsubsection{Circuito de polarización fija}
???
\subsubsection{Circuito de auto-polarizado}
???
\section{Transistores PNP}
???
\subsection{Curvas características del transistor PNP}
???
\subsection{Modelos del transistor PNP en corriente continua}
???
\section{Amplificadores con transistores}
???
\subsection{Configuración emisor común}
???
\subsection{Amplificador base común}
???
\subsection{Amplificador colector común}
???
\subsection{Amplificador emisor común con $R_E$}
???
\section{Máxima señal de salida sin distorsión. Rectas de carga}
???
\subsection{Amplificador simple emisor común}
???
\subsubsection{Rectas de carga}
???
\subsubsection{Distorsión}
???
\subsubsection{Magnitud de la máxima señal de salida sin distorsión}
???
\subsubsection{Punto $Q$ óptimo}
???
\subsection{Amplificador con carga acoplada capacitivamente}
???
\subsubsection{Rectas de carga en $DC$}
???
\subsubsection{Recta de carga en alterna}
???
\subsubsection{Magnitud de la máxima señal de salida}
???
\subsection{Otras configuraciones}
???
\subsubsection{Amplificador emisor común con $R_s$}
???
\subsubsection{Amplificador base común}
???
\subsubsection{Amplificador emisor común con $CE$ y carga acoplada capacitivamente}
???
\subsection{Diseño de amplificadores con transistores}
???